% ======================================================================
%  Eidesstattliche Erklaerung
%
%  Der Wortlaut ist pruefungsordnungsabhaengig -- vor der Abgabe gegen die
%  Vorlage der eigenen Hochschule pruefen und gegebenenfalls ersetzen. Der
%  Absatz zu KI-Werkzeugen ist inzwischen vielerorts vorgeschrieben; wo er
%  nicht verlangt ist, kann er entfallen.
% ======================================================================

\cleardoublepage
\chapter*{Eidesstattliche Erklärung}
\addcontentsline{toc}{chapter}{Eidesstattliche Erklärung}
\markboth{Eidesstattliche Erklärung}{Eidesstattliche Erklärung}

Ich versichere an Eides statt, dass ich die vorliegende Bachelorarbeit mit dem
Titel

\begin{quote}
\enquote{Datenaugmentierung für die CNN-basierte Detektion von
Rohrprofil-Stirnflächen an gestapelten Bauzaunfeldern}
\end{quote}

\noindent selbstständig und ohne fremde Hilfe verfasst und keine anderen als
die angegebenen Quellen und Hilfsmittel benutzt habe. Alle Stellen, die
wörtlich oder sinngemäß aus Veröffentlichungen entnommen sind, habe ich als
solche kenntlich gemacht. Die Arbeit hat in gleicher oder ähnlicher Form noch
keiner Prüfungsbehörde vorgelegen.

\platzhalter{Sofern von der Prüfungsordnung verlangt: Angabe zu Art und Umfang
der Nutzung KI-gestützter Werkzeuge.}

\vspace{2.5cm}

\noindent
\begin{tabular}{@{}p{0.45\textwidth}p{0.45\textwidth}@{}}
\hrulefill & \hrulefill \\
\platzhalter{Ort}, \platzhalter{Datum} & Moritz Stötzner \\
\end{tabular}
